\chapter*{术语约定}
\addcontentsline{toc}{chapter}{术语约定}

\begin{longtable}{>{\raggedright\arraybackslash}p{4.0cm}>{\raggedright\arraybackslash}p{9.6cm}}
\caption{本文使用的核心术语}\label{tab:terminology}\\
\toprule
\textbf{术语} & \textbf{含义}\\
\midrule
主体（Subject） & 跨模型、会话和运行进程保持逻辑身份的长期主体。\\
主体标识（SubjectId） & 在限界上下文之间保持一致的主体逻辑标识。\\
认知权威（Cognitive Authority） & 对持久认知对象、修订、生命周期和写入权限具有决定作用的数据与领域合同。\\
认知运行时（Cognitive Runtime） & 组织当前活动认知、工作上下文、注意分配和调用方工作集的可恢复运行状态。\\
服务投影（Serving Projection） & 从权威数据构建并可独立重建的索引、向量、拓扑和其他读优化视图。\\
观察事件（ObservationOccurrence） & 主体在特定时间实际接触某份材料或事件的一次记录。\\
证据引用（EvidenceRef） & 指向观察事件、来源区域或派生表示的可追溯引用。\\
记忆（Memory） & 由记忆域接受并长期保存的主体认知。\\
认知图式（CognitiveSchema） & 主体从多段经历和认知中形成的、带有适用范围与反例的结构性认知。\\
社会认知（Social Cognition） & 依赖具体个人、关系、群体、频道或社区语境的长期社会性认知。\\
假设（Hypothesis） & 需要由目标领域所有者核验后才能形成权威认知的候选主张。\\
工作上下文（WorkContext） & 可暂停、恢复或结束的任务或问题上下文。\\
\bottomrule
\end{longtable}
\clearpage
